% =============================================================================
% Appendix E  The Structured Exact Encoding of the Differentiation Matrix
% Cited from sec:standard-factorisation and sec:standard-cost.
% Plan: ThesisPlanning.md §6 (E); supporting contribution S2 (claims C-01, C-02).
%
% CONTENT  The proof of prop:G-factorisation; Lambda, U_odd (subtraction with a
%          borrow flag) and R_0 as circuits; the resource derivation behind
%          prop:standard-cost; alpha / ||G||_2 across n.
% SOURCES  Planning.md §3.7 ("Standard form, structured exact"), App. A-4;
%          pihm/src/pihm/circuits/standard.py (structured_derivative_encoding),
%          arithmetic.py (register_subtractor); pihm/docs/PILLAR1.md.
% =============================================================================

\subsection{Proof of \texorpdfstring{\Cref{prop:G-factorisation}}{the factorisation}}
\label{app:proof-G-factorisation}
% LAND: entry by entry -- (R_0 (2 U_odd) Lambda)_{kp} = r_k * 2[p > k, p - k odd]
%   * p, with r_0 = 1/sqrt2 and r_k = 1 otherwise, which is eq:pihm-G-entries.
%   Then Lambda's Z-string form (eq:lambda-lcu) from p = sum_j 2^j (1 - Z_j)/2.

\subsection{The arithmetic SELECT}
\label{app:uodd-select-detail}
% LAND: U_odd as the average of the N/2 truncated shifts by odd offsets;
%   PREPARE (X on the offset register's lowest bit, H on the rest); SELECT as
%   the in-place subtraction p <- p - d (Cuccaro ripple-carry,
%   cuccaro2004adder) with a borrow flag post-selected on |0>, so every wrapped
%   term vanishes exactly. Toffoli count O(n). The circuit is fig:uodd-select.

\subsection{Proof of \texorpdfstring{\Cref{prop:standard-cost}}{the cost of the structured encoding}}
\label{app:proof-standard-cost}
% LAND: alpha multiplies under products: alpha(R_0) = 1, alpha(2 U_odd) = N,
%   alpha(Lambda) = N - 1, so alpha_G = N(N - 1). Gate count O(n^2) and ancilla
%   n + O(log n), with the Omega(n) floor from the N/2-term LCU. Products reuse
%   ancilla with one failure flag per product (implementation issue I-03: until
%   built, products allocate fresh ancilla -- say which the reported counts use).

\subsection{Subnormalisation across \texorpdfstring{$n$}{n}}
\label{app:alpha-G}
% The ratios below are the planning probe's (Planning.md App. A-4; the
%   factorisation's error is <= 2.3e-13 throughout); regenerate from results-v1.

\begin{table}[htbp]
    \centering
    \caption{Subnormalisation of the structured and published encodings of $\G$, relative to $\|\G\|_2$.}
    \label{tab:alpha-G-full}
    \small
    \begin{tabular}{@{}r r l l@{}}
        \toprule
        $n$ & $\|\G\|_2$ & Structured & Published \\
        \midrule
        2 & \res[3,fixed]{derivative/n2}{norm} & \res[4]{derivative/n2}{structured_ratio} & \res[3]{derivative/n2}{published_ratio} \\
        3 & \res[3,fixed]{derivative/n3}{norm} & \res[4]{derivative/n3}{structured_ratio} & \res[3]{derivative/n3}{published_ratio} \\
        4 & \res[3,fixed]{derivative/n4}{norm} & \res[4]{derivative/n4}{structured_ratio} & \res[3,sci]{derivative/n4}{published_ratio} \\
        5 & \res[3,fixed]{derivative/n5}{norm} & \res[4]{derivative/n5}{structured_ratio} & \res[3,sci]{derivative/n5}{published_ratio} \\
        6 & \res[3,fixed]{derivative/n6}{norm} & \res[4]{derivative/n6}{structured_ratio} & \res[3,sci]{derivative/n6}{published_ratio} \\
        7 & \res[3,fixed]{derivative/n7}{norm} & \res[4]{derivative/n7}{structured_ratio} & \res[3,sci]{derivative/n7}{published_ratio} \\
        8 & \res[3,fixed]{derivative/n8}{norm} & \res[4]{derivative/n8}{structured_ratio} & \res[3,sci]{derivative/n8}{published_ratio} \\
        9 & \res[3,fixed]{derivative/n9}{norm} & \res[4]{derivative/n9}{structured_ratio} & \res[3,sci]{derivative/n9}{published_ratio} \\
        10 & \res[3,fixed]{derivative/n10}{norm} & \res[4]{derivative/n10}{structured_ratio} & \res[3,sci]{derivative/n10}{published_ratio} \\
        \bottomrule
    \end{tabular}
\end{table}
